\documentclass[UTF8,11pt]{ctexart}
\usepackage[a4paper,margin=2.3cm]{geometry}
\usepackage{amsmath,amssymb,amsthm}
\usepackage{booktabs}
\usepackage{graphicx}
\usepackage{float}
\usepackage[numbers]{natbib}
\usepackage{hyperref}
\usepackage{xcolor}
\usepackage{enumitem}
\usepackage{array}
\usepackage{longtable}
\usepackage{caption}
\graphicspath{{figures/}{./}}
\hypersetup{colorlinks=true,linkcolor=blue,citecolor=blue,urlcolor=blue}
\captionsetup{font=small}

\newtheorem{proposition}{命题}
\newtheorem{remark}{评注}
\newtheorem{assumption}{假设}
\newtheorem{definition}{定义}
\newcommand{\qos}{q}
\newcommand{\T}{\mathcal{T}}
\newcommand{\M}{\mathcal{M}}
\newcommand{\D}{D}
\newcommand{\w}{\boldsymbol{w}}
\newcommand{\p}{\boldsymbol{p}}
\newcommand{\g}{\boldsymbol{g}}

\title{\textbf{算力固定下的推理服务分时动态定价：\\
异质厂家竞争、用户负载迁移与QoS保护机制}}
\author{匿名作者}
\date{\today}

\begin{document}
\maketitle

\begin{abstract}
大模型推理服务的GPU算力在短期内基本固定，而请求负载具有显著的时段性峰谷差：高峰时段过载导致拥塞与服务质量（QoS）退化，低谷时段算力闲置。本文研究一个现实问题：在算力固定的条件下，厂家能否通过\textbf{分时动态定价}（高峰提价、低谷降价）诱导时间弹性负载从拥塞高峰迁移到空闲低谷，从而在不扩容的前提下保护服务质量。我们建立一个异质双厂家竞争市场：两家部署同款开源模型（如GLM）、服务近似同质的厂家各自拥有固定算力，一个可在两家之间路由的中间商，以及两类用户——时间刚性用户（交互式、实时，迁移成本高）与时间弹性用户（批量、离线，对低谷低价敏感）；用户还可经厂家直连API或外部退出选项离开。我们用低维形状参数化将每个厂家的分时定价压缩为少数标量，并以\textbf{虚拟博弈（fictitious play）}求解厂家间的Nash竞争——这一选择解决了朴素最优响应在拥塞区间出现的极限环振荡问题，使均衡求解可复现（非拥塞区间跨随机种子相对标准差为$4\times10^{-4}$）。

实验给出一个经过对抗性鲁棒性检验的核心结论：\textbf{分时动态定价在市场拥塞时是可靠的QoS保护工具，但对厂家利润的影响近似中性，而非牟利手段}。在拥塞区间，统一定价使最低QoS跌至$0.756$、峰值利用率达$0.782$；动态定价将峰值利用率削至$0.66$--$0.71$、最低QoS恢复到$0.97$--$0.99$——这一QoS改善在粗、细两种网格分辨率下方向一致，是稳健的。但系统利润的变化在粗网格下为$+9.3\%$、在细网格下为$-1.9\%$，符号随求解精度翻转（且细网格求解未收敛），表明真实利润效应接近于零。我们从机制上解释这一"利润打平"：在算力与需求总量固定时，分时定价主要\textbf{重新分配剩余}而非创造剩余——高峰提价的增益被驱离的需求抵消，低谷填谷的增量被薄利抵消，而QoS改善带来的剩余在竞争与退出选项的约束下流向用户而非厂家利润。这一发现与"动态定价是厂家趁拥塞牟利"的通常预期相反，具有监管与运营含义：分时定价的正当理由是保障SLA与服务质量，而非提升收入。本文结论受限于合成参数标定，具体数值仅支撑机制方向，不作生产预测。
\end{abstract}

\noindent\textbf{关键词：}推理服务；动态定价；削峰填谷；异质厂家竞争；负载迁移；QoS保护；虚拟博弈；剩余重分配

% ============================================================
\section{引言}

\subsection{研究背景与问题}

大模型推理服务已经从离线批处理转向持续运行的在线服务。交互式聊天、代码生成、自动代理、批量评测和定时工作流在一天内呈现不同的到达规律，使GPU推理集群面临显著的峰谷差\cite{crankshaw2017clipper,gujarati2020clockwork,yu2022orca,kwon2023pagedattention,agrawal2024sarathiserve}。当利用率接近容量边界时，首token延迟、排队时间和有效完成率快速恶化，服务质量（QoS）退化直接影响可结算服务量和用户满意度。

与可以快速调度发电或长期扩容的系统不同，推理服务的GPU算力在短期内\textbf{基本固定}——采购和部署有显著的前置时间和沉没成本，且低谷闲置的算力照样产生持有成本。这造成一个核心张力：高峰时段算力不足导致拥塞与QoS退化，低谷时段算力闲置造成浪费。一个自然的问题是：能否不依赖扩容，而是用\textbf{价格信号}把负载在时间上重新分配？

本文研究这一问题的一个具体形式：在算力固定的条件下，厂家通过\textbf{分时动态定价}（高峰提价、低谷降价），能否诱导对时间不敏感的负载从拥塞高峰迁移到空闲低谷，从而保护高峰期的服务质量。这一思路与电力市场的实时定价和需求响应同源\cite{schweppe1988spot,borenstein2005long,albadi2008summary,faruqui2010household,siano2014demand}，但推理服务的可计算约束不同：服务质量由GPU局部利用率、延迟和完成率决定，而非发电约束与网络潮流。

\subsection{为什么需要竞争与异质性}

现实推理服务市场并非单一厂家垄断。多个厂家可以部署同一款开源模型（例如GLM系列），向用户提供\textbf{近似同质}的服务，竞争主要落在价格与算力充裕度（QoS）上，而非模型能力差异。这一观察给竞争提供了真实的微观基础，也使本文与"单平台垄断定价"的简化区分开来。

此外，厂家的算力规模通常\textbf{异质}：老牌厂家算力充裕，新入场厂家算力紧张。算力异质会改变动态定价的策略空间——算力紧张的厂家更容易被路由来的流量打到拥塞。本文因此设定两家算力不同但服务同质的厂家，并保留一个可在两家之间路由的中间商（如OpenRouter类API聚合商），以及可绕过中间商直连厂家API或退出市场的用户。

\subsection{贡献}

\begin{enumerate}[leftmargin=2em]
  \item \textbf{算力固定下的削峰填谷动态定价模型。} 建立异质双厂家竞争、中间商路由、两类用户（时间刚性/弹性）加外部退出选项的市场模型，将厂家的固定算力闲置成本、QoS拥塞退化和用户负载迁移纳入同一可计算框架。
  \item \textbf{可复现的均衡求解。} 用低维形状参数化加虚拟博弈（fictitious play）求解厂家间Nash竞争。我们识别并解决了朴素最优响应在拥塞区间的极限环振荡问题——这是同类多层定价博弈反复出现的收敛障碍；本文的求解在非拥塞区间达到跨种子相对标准差$4\times10^{-4}$，在拥塞区间以regret判据单调收敛。
  \item \textbf{经鲁棒性检验的QoS保护机制发现。} 经粗、细网格的对抗性鲁棒性检验，给出一个诚实结论：分时动态定价在拥塞时\textbf{可靠保护QoS（削峰、恢复完成率）但对厂家利润近似中性}。我们从剩余重分配的角度解释这一"利润打平"，并指出其与"动态定价牟利"预期相反的监管与运营含义。
\end{enumerate}

% ============================================================
\section{模型}
\label{sec:model}

\subsection{市场结构与时序}

一天划分为$T$个时段$\T=\{1,\ldots,T\}$，每时段时长$h$。市场含：
\begin{itemize}[leftmargin=2em]
  \item \textbf{两个异质厂家} $\M=\{A,B\}$，固定算力$G_A>G_B$，部署同款开源模型（服务同质）。厂家$m$选择分时批发价$w_{m,t}$（卖给中间商）和分时直连零售价$p^D_{m,t}$（卖给用户）。
  \item \textbf{一个中间商}，观察批发价后选择对用户的分时零售价$p_t$和把每时段流量分配给两厂家的路由权重$r_{m,t}$（$\sum_m r_{m,t}=1$）。
  \item \textbf{两类用户}，按logit随机效用在"中间商通道、厂家A直连、厂家B直连、退出"之间选择。
\end{itemize}
时序为：厂家同时出价（Nash竞争）$\to$ 中间商最优响应 $\to$ 用户选择与QoS拥塞固定点。

\subsection{两类用户与外部退出选项}

设用户类型$\theta\in\{R,E\}$，$R$为时间刚性、$E$为时间弹性。类型$\theta$对通道$k$（中间商/厂家A/厂家B）在时段$t$的确定性效用为
\begin{equation}
  z^\theta_{k,t}=\log b_k+\log a_t-\alpha_\theta\big(\text{price}_{k,t}-p_0\big)
  -c_s^\theta(1-\pi^\theta_t)-\gamma(1-q_{k,t}),
  \label{eq:utility}
\end{equation}
其中$b_k$为通道品牌/便利性，$a_t$为时段偏好，$\pi^\theta_t$为类型$\theta$的原生时段分布，$q_{k,t}$为通道QoS。类型差异完全由参数承载：刚性用户价格敏感度$\alpha_R$低、迁移成本$c_s^R$高、原生时段集中于高峰；弹性用户$\alpha_E$高、$c_s^E$低、原生时段分散。

引入确定效用为$u_0$的\textbf{no-purchase外部退出选项}进入logit归一化。类型$\theta$选择通道$k$时段$t$的份额与退出概率为
\begin{equation}
  s^\theta_{k,t}=\frac{\exp(z^\theta_{k,t})}{\exp(u_0)+\sum_{k',t'}\exp(z^\theta_{k',t'})},
  \qquad
  s^\theta_0=\frac{\exp(u_0)}{\exp(u_0)+\sum_{k',t'}\exp(z^\theta_{k',t'})}.
  \label{eq:share}
\end{equation}
内部份额之和$<1$，残余质量$s^\theta_0$为真实退出概率，不被重新分配回内部选项。这保证了价格抬高时需求真正流失，而非被归一化掩盖。

\subsection{需求、负载与QoS}

类型$\theta$对通道$k$的需求由弹性项和刚性/重放项构成，二者均\textbf{按渠道自身份额}分配（使每时段刚性总量与渠道数无关）：
\begin{equation}
  \D^\theta_{k,t}=\bar F_\theta\,m(\p)\,s^\theta_{k,t}
  +\bar R^\theta_t\,\sigma[\beta(v_r-\text{price}_{k,t})]\,s^\theta_{k,t}.
  \label{eq:demand}
\end{equation}
总需求为两类人口加权之和，权重$(\pi_R,\pi_E)$。厂家$m$在时段$t$的总负载为直连需求加中间商按路由分来的需求：
\begin{equation}
  L_{m,t}=\D^D_{m,t}+r_{m,t}\,\D^I_t,\qquad u_{m,t}=L_{m,t}/G_m,\qquad q_{m,t}=\qos(u_{m,t}).
  \label{eq:load}
\end{equation}
QoS退化函数$\qos(u)$在利用率超过阈值$\bar u$后下降。算力异质$G_B<G_A$使同等负载下小厂$B$更易越过阈值而拥塞。本文主用平滑sigmoid退化形态（处处可微，利于均衡求解），并对阈值型、线性、根号形态做鲁棒性检验。

\subsection{利润：含固定算力闲置成本}

厂家$m$的利润为批发收入加直连收入，减去\textbf{按固定算力计的闲置成本}与QoS退化成本：
\begin{equation}
  \Pi_m=h\!\sum_t\! w_{m,t}r_{m,t}\D^I_t q_{m,t}
  +h\!\sum_t\! p^D_{m,t}\D^D_{m,t}q_{m,t}
  -h\!\sum_t\! c_g G_m
  -h\!\sum_t\! c_q(1-q_{m,t})L_{m,t}.
  \label{eq:firm_profit}
\end{equation}
闲置成本$c_g G_m$按固定算力计——无论是否用满都付，这是厂家有动机降低谷价以填谷的根源，也是本文与"容量为常数、无闲置成本"模型的关键区别。中间商利润为零售收入减批发成本减QoS退化成本。系统利润为两厂家利润加中间商利润。


% ============================================================
\section{均衡求解}
\label{sec:solver}

\subsection{低维形状参数化}

直接让每个厂家在$2T$维（分时批发价加直连价）连续空间中优化，会使Nash求解高维、非凸、不可复现——这是同类多层定价博弈反复出现的收敛障碍。本文将厂家的分时定价\textbf{参数化为少数标量形状}：
\begin{equation}
  w_{m,t}=\bar w_m\big(1+\delta_m\,\widehat{\ell}_t\big),\qquad
  p^D_{m,t}=\bar p^D_m\big(1+\delta^D_m\,\widehat{\ell}_t\big),
  \label{eq:shape}
\end{equation}
其中$\widehat{\ell}_t$是零均值的标准化负载形状（由刚性需求基线归一化得到，外生已知），每个厂家的策略压缩为四个标量$(\bar w_m,\delta_m,\bar p^D_m,\delta^D_m)$。动态强度$\delta_m=0$对应统一定价；$\delta_m>0$对应高峰提价。这一约束以"形状族内最优"换取可复现性：本文不声称$2T$维连续空间的全局最优，只声称该形状族内的均衡。

\subsection{虚拟博弈与极限环的解决}

厂家的最优响应通过对四标量的\textbf{粗网格搜索}计算（对非光滑QoS目标鲁棒，成本可预测、可并行）；中间商响应同样用网格搜索，每次评估内嵌QoS与路由的\textbf{联合固定点}（一次阻尼迭代，而非嵌套求解，避免拥塞下的成本爆炸）。

朴素的最优响应迭代（厂家轮流对对手\emph{当前}策略最优响应）在拥塞区间出现\textbf{极限环振荡}：策略在两组配置间循环往复，永不收敛。这与寡头定价博弈的已知现象一致。本文改用\textbf{虚拟博弈（fictitious play）}：每个厂家对对手的\textbf{运行平均}策略最优响应，历史平均天然阻尼振荡。收敛判据采用\textbf{regret}（厂家相对时间平均策略的单边最大偏离收益）而非参数漂移——后者在求解器分辨率下会误报。

\begin{remark}[收敛性，含诚实边界]
非拥塞区间，QoS恒为1、固定点单步收敛，网格求解完全确定性，跨随机种子系统利润相对标准差为$4.31\times10^{-4}$，收敛性完全可靠。拥塞区间则需谨慎：在\textbf{粗网格}下虚拟博弈的regret单调下降（$133.6\to63.7\to\cdots\to4.98$，22轮达到regret$<5$）；但在\textbf{细网格}下regret在约第10轮触底于$\approx14$后不降反升、40轮后漂移到$\approx22$，\textbf{未能达到收敛判据}。这说明粗网格的"regret$=4.98$收敛"部分是粗粒度掩盖了可盈利偏离的假象——更细的偏离候选暴露出该拥塞博弈在本文求解预算内\textbf{并非可靠收敛}。因此本文对拥塞区间只作分辨率稳健的定性结论（见第~\ref{sec:results} 节），不依赖任何依赖完全收敛的定量声明。
\end{remark}

% ============================================================
\section{实验结果}
\label{sec:results}

\subsection{设置}

$T=8$时段、$h=3$小时。基准算力$G_A=1500,G_B=600$；拥塞情形将算力收紧至$G_A=300,G_B=120$以激活QoS退化。两类用户人口占比基准$(\pi_R,\pi_E)=(0.6,0.4)$；拥塞情形取$(0.4,0.6)$以加大高峰压力。所有经济参数为合成标定，具体数值仅支撑机制方向，不作生产预测。完整参数见补充材料。

\subsection{非拥塞区间：负载迁移稳健，利润中性}

基准算力下市场不拥塞（统一定价峰值利用率仅$0.213$），QoS恒为1。表~\ref{tab:uncongested} 给出统一与动态定价对比。

\begin{table}[H]
\centering
\caption{非拥塞区间：统一定价 vs 动态定价（确定性网格求解）}
\label{tab:uncongested}
\begin{tabular}{lccc}
\toprule
指标 & 统一定价 & 动态定价 \\
\midrule
系统利润 & 1193.3 & 1139.4 \\
峰值利用率 & 0.213 & 0.261 \\
弹性需求时段重心位移 & \multicolumn{2}{c}{$+0.177$} \\
刚性需求时段重心位移 & \multicolumn{2}{c}{$+0.077$} \\
\bottomrule
\end{tabular}
\end{table}

\textbf{负载迁移（P1）稳健}：动态定价下弹性需求的时段重心后移$+0.177$，刚性需求仅$+0.077$——弹性负载迁移强度约为刚性的$2.3$倍，符合"动态定价主要驱动时间弹性负载迁移"的机制预期。但在不拥塞的市场，QoS本就为1、无可保护，动态定价的利润效应中性到微负（弹性占比扫描下增益为$0\%,-0.32\%,-4.51\%,-3.76\%,-3.00\%$）。这说明动态定价的价值前提是\textbf{市场确实拥塞}。

\subsection{拥塞区间：QoS保护稳健，利润打平}

收紧算力后，统一定价使市场真正拥塞。表~\ref{tab:congested} 给出虚拟博弈求解的结果，并列粗、细两种网格分辨率以检验鲁棒性。

\begin{table}[H]
\centering
\caption{拥塞区间：统一 vs 动态定价（虚拟博弈均衡），含网格鲁棒性检验}
\label{tab:congested}
\begin{tabular}{lccc}
\toprule
指标 & 统一定价 & 动态（粗网格） & 动态（细网格） \\
\midrule
峰值利用率 & 0.782 & 0.706 & 0.666 \\
最低QoS & 0.756 & 0.968 & 0.990 \\
系统利润 & 1783 & 1949 & 1749 \\
利润变化 & --- & $+9.3\%$ & $-1.9\%$ \\
\bottomrule
\end{tabular}
\end{table}

\textbf{QoS保护（P2、P3）稳健、分辨率无关。} 统一定价使最低QoS跌至$0.756$、峰值利用率达$0.782$。动态定价将峰值利用率削至$0.66$--$0.71$、最低QoS恢复到$0.97$--$0.99$。削峰与QoS保护在粗、细两种网格下方向一致，是稳健结论。

\textbf{利润效应近似中性、不稳健。} 系统利润变化在粗网格下为$+9.3\%$，在细网格下翻转为$-1.9\%$。利润变化的\textbf{符号本身随求解精度翻转}，且细网格虚拟博弈本身未收敛（regret漂移至$\approx22$，见第~\ref{sec:solver} 节评注），故真实利润效应接近于零、被求解误差淹没。因此本文只声称动态定价利润\textbf{近似中性}，不声称提升利润。

\textbf{被撤回的假设。} 我们曾假设小算力厂家更依赖动态定价（动态强度$\delta_B>\delta_A$）。粗网格给出$\delta_A=0.383\gg\delta_B=0.023$（大厂主导），细网格给出$\delta_A=0.059,\delta_B=0.278$（小厂主导）——方向随网格翻转。该假设在任一方向上都不稳健，应予撤回。同样，负载迁移非对称性（P1）在非拥塞区稳健，但在拥塞均衡处抹平（细网格弹性/刚性重心$4.23$对$4.31$，差异微弱）。

% ============================================================
\section{讨论：为什么利润打平}
\label{sec:discussion}

拥塞区间的核心现象是QoS大幅改善（$0.756\to0.97$以上）而利润近似不变。这并非数值巧合，而是有经济机制的结果。

\subsection{剩余重分配而非创造}

关键前提是\textbf{算力固定、需求总量基本固定}。动态定价只在时间维度上重排需求，不创造新算力或新需求。其对厂家利润的影响分解为方向相反的两组力量。正向：高峰提价提高单位毛利；低谷降价利用闲置算力产生额外结算量；削峰使QoS回升、有效结算量增加。负向：高峰提价驱离价格敏感需求（流向对手、直连或退出）；低谷降价摊薄单位毛利；被挤出的弹性用户部分净退出。高峰"提价多赚"被"驱离需求少赚"抵消，低谷"填谷多做量"被"低价薄利"抵消，厂家总利润近似不变。这与价格歧视的标准结论一致：供给与需求总量固定时，时间价格歧视主要重新分配剩余，不必然放大厂家利润。

\subsection{削峰省下的损失为何没变成利润}

QoS从$0.756$回升到$0.97$以上、少折损大量请求，直觉上应是纯增益。两个原因使其未落到厂家利润：第一，\textbf{削峰是靠驱离高峰需求实现的}——把利用率从$0.78$压到$0.66$必须挤走部分高峰用户，QoS回升多结算的量大致等于被挤走需求的损失，保QoS的代价恰是丢需求。第二，\textbf{竞争与退出选项封住涨价空间}——异质双厂家竞争加直连/退出选项，使任一厂家无法靠提价独吞QoS改善的剩余；该剩余更多流向用户（体验改善、价格被竞争压制），而非厂家利润。

\subsection{政策与运营含义}

这一发现与"动态定价/surge pricing是厂家趁拥塞牟利的工具"的通常预期相反。在竞争、算力固定、用户可退出的推理服务市场，分时定价\textbf{主要是QoS保护工具，厂家并不因此显著多赚}。监管含义：监管者不必将分时定价视为牟利手段，其剩余主要流向服务质量与用户体验。运营含义：运营者采用分时定价的正当理由应是保障SLA与服务质量，而非提升收入——把它当收入工具会落空。

% ============================================================
\section{局限}
\label{sec:limitations}

第一，全部经济参数（价格敏感度、迁移成本、容量成本、弹性人口占比、退出效用）为合成标定，无真实请求轨迹校准；具体百分比不作生产预测，仅支撑机制方向。第二，厂家策略限定在低维形状族内，不声称连续策略空间的全局最优；更重要的是，\textbf{拥塞区间的虚拟博弈在细网格下并未收敛}——regret在约第10轮触底于$\approx14$后漂移至40轮的$\approx22$，远未达到$<5$判据，故拥塞细网格的利润与策略数（含$-1.9\%$、$\delta$值）均为未收敛快照，仅用于说明"利润方向不稳健"。粗网格的"regret$=4.98$收敛"则部分是粗粒度掩盖可盈利偏离的假象。这意味着拥塞区间的\textbf{定量}结论（利润、哪个厂家更动态）尚不可靠，只有\textbf{定性}的QoS保护方向（削峰、最低QoS回升）在粗细网格下一致，可作为本文的稳健主张。彻底解决拥塞均衡的可靠求解（更强的平均化、变分不等式重构或更细的收敛预算）是必要的后续工作。第三，本文为两厂家、单中间商结构；更多厂家、多中间商、长期容量合约与模型路由的复杂交互留待后续。第四，QoS退化用合成sigmoid曲线，真实后端的完成率/延迟曲线需用受控测量或生产轨迹替换重标。第五，利润打平的结论建立在"算力与需求总量固定"前提上；若需求总量本身对价格高度弹性、或算力可短期调整，结论可能改变。

% ============================================================
\section{结论}
\label{sec:conclusion}

本文研究算力固定下推理服务的分时动态定价能否保护服务质量。我们建立异质双厂家竞争、中间商路由、两类用户加退出选项的市场模型，用低维形状参数化加虚拟博弈求解厂家Nash竞争，解决了朴素最优响应在拥塞区间的极限环振荡，使均衡求解可复现。经粗、细网格的对抗性鲁棒性检验，核心结论是：\textbf{分时动态定价在市场拥塞时是可靠的QoS保护工具}（削峰、把最低QoS从$0.756$恢复到$0.97$以上，分辨率无关），\textbf{但对厂家利润近似中性}（利润变化符号随求解精度翻转）。我们从剩余重分配的角度解释这一利润打平：算力与需求固定时，分时定价重排而非放大剩余，QoS改善的剩余在竞争与退出约束下流向用户。这一发现与"动态定价牟利"预期相反，其监管与运营含义是：分时定价的正当理由是保障服务质量，而非提升收入。未来工作包括真实轨迹校准、拥塞区间的完全收敛与置信区间、以及多厂家多中间商扩展。

\section*{代码与工件可用性}
削峰填谷仿真模块为\path{pricing_sim/peak_shaving_config.py}、\path{pricing_sim/peak_shaving_market.py}、\path{pricing_sim/peak_shaving_equilibrium.py}；实验入口为\path{experiments/run_peak_shaving_experiments.py}（非拥塞主实验）、\path{experiments/run_peak_shaving_congested_fp.py}（拥塞虚拟博弈）、\path{experiments/run_fp_dynamic_converge.py}（regret收敛与断点续跑）、\path{experiments/run_peak_shaving_robustness.py}（QoS形态、退出效用、跨种子）。机器可读工件在\path{artifacts/peak_shaving/20260618/}，每个正文数字对应一条落盘JSON记录。

\small
\bibliographystyle{unsrtnat}
\bibliography{verified_refs}

\end{document}
